\chapter{Epilogue: The Wider Frame}
\label{ch:epilogue}

The practical questions of the previous chapters sit in a wider context. They
connect the systems we build to questions about evolution, the biosphere and
technological change, which ecologists have a professional reason to consider.

\section{Evolution and the biosphere}

\textcite{muller2026evolvableai} ask when AI systems whose components and
learning rules can themselves evolve by natural selection become possible, and
argue the scenario is underrated. Treating deployed AI as an evolving population
is a move ecologists are well placed to judge. \textcite{korecki2025biosphericai}
argues that AI ethics centres on humans, at most on sentient beings, and
proposes the biosphere as the design target instead. Both connect to Vernadsky's
idea of a transition from biosphere to \term{noosphere}, a sphere of thought
\parencite{vernadsky1938noosphere}. \textcite{vidal2024noosphere} revisit it as
a major evolutionary transition, and \textcite{lovelock2019novacene} pushed it
further, arguing that a new age of AI beings is beginning. \textcite{vonneumann1955survive} asked the version that has aged
best, in 1955: whether the world has become too small, and its political units
too unstable, for the technology being built in it.

Benedikt Hartl, at the Allen Discovery Center at Tufts, approaches the same
transition from developmental biology. In his perspective, the genome is a compact
library of past survival lessons that cells consult as they develop.
Language models, he suggests, could come to play that part in cultural
evolution, as a generative memory of collective knowledge that gives
individuals skills they never learned \parencite{hartl2026moreisdifferent}. He
reads this as the start of another major transition, and asks whether people,
machines, institutions, ecosystems and the biosphere could form one higher
level of organisation. He also notes how far we are from it, since a
connected society still fails to act together on threats such as biodiversity
loss: ``globally connected, but not yet globally coherent''. This is an
argument, not a result.

The same centre also works the comparison the other way, using a method from
AI to model development. \textcite{hartl2026growthform} applied low-rank adaptation, a
technique for fine-tuning language models, to neural cellular automata, in
which cells grow a shape from a single seed by following a shared local rule.
Rank-one changes to that rule reproduced the horizontal and vertical stretching
of D'Arcy Thompson's grid transformations (\Cref{ch:landscape}), and carried
over without retraining to other shapes grown from the same rule. The shapes
are emoji on a two-dimensional grid, and the authors note that the model's
weights cannot yet be assigned to genes or signals.

\section{An inherited archive}

David Krakauer's essay on the fossil logic of AI, listed in \Cref{app:tutorials},
adds the history of accumulated knowledge to this frame. He compares extracting
useful work from an inherited archive with extracting energy from fossil fuels.
The difference matters: reading does not consume knowledge, but returning
unchecked generated claims to the archive can damage its reliability. This is
an interpretive analogy, not a measured ecological result. Machine learning
offers a measured counterpart. Models trained repeatedly on generated text lose
the rare parts of the original distribution first \parencite{shumailov2024collapse};
keeping enough real data in the mix can avoid this \parencite{seddik2024collapse}.
For biodiversity, I see the same risk in generated occurrence records, trait
values and phenotype assignments: once deposited in public databases, they
become other people's evidence and training data (\Cref{sec:outputasdata}).
I take the practical implication to be continued investment in collections,
field observations and careful curation. These activities renew the evidence
on which later models and researchers depend.

\section{Consciousness}

Consciousness belongs in this frame too. Like evolvable AI and biospheric
ethics, it turns on how built systems relate to living ones.
\textcite{chandaria2026consciousness}, introduced in \Cref{ch:landscape},
organise competing theories around five things to examine:
behaviour, computation, physical causal organisation, bodily self-maintenance,
and interaction between an organism and its environment. They combine evidence
with explicit assumptions about which of these matters for consciousness.
Their numerical examples illustrate how conclusions change with those
assumptions; they are not measured probabilities that a particular model is
conscious. For readers coming from biology, the framework makes the disagreement
clear: do the functions suffice, or do the living processes that realise them
matter too?

An exchange in \emph{Behavioral and Brain Sciences} makes that disagreement
concrete. \textcite{seth2026naturalism} argues that consciousness may depend on
the processes by which living systems maintain themselves. He regards
consciousness in conventional digital AI as unlikely, while leaving open
artificial systems that reproduce relevant properties of brains or life.
He does not claim that all life is conscious, or that conscious life must be
carbon-based. In a later \emph{Nature} comment, \textcite{seth2026threereasons}
give three reasons for scepticism: intelligence is not consciousness, brain
function cannot be separated from living tissue, and several theories of
consciousness are not computational.

\textcite{richards2026computational} agree that living processes matter but
argue that prediction, bodily regulation and self-production are themselves
computational. They therefore consider artificial consciousness likely and
suggest that some current systems might already have limited experience.
They present this as a hypothesis, not an established observation.
The disagreement concerns what would be sufficient for experience; neither
position makes fluent conversation a decisive test
\parencite{seth2026naturalism,richards2026computational}.

The framework of \textcite{chandaria2026consciousness} also raises an
ethical uncertainty. Mistakenly attributing experience to a system can
divert concern from humans and other animals; overlooking experience, if it
exists, could conceal suffering. The authors explicitly leave disagreement
about consciousness unresolved. An ecological accuracy score cannot resolve it
either. My judgement is to keep these questions visible while assessing
research systems through their observable effects on data, people and the
biosphere.
